\makeatletter
\@ifundefined{keepfig}{\newenvironment{keepfig}{\par\edef\keepfig@pi{\the\parindent}\noindent
  \begin{minipage}{\linewidth}\parindent=\keepfig@pi\relax}{\end{minipage}\par}}{}
\makeatother
\section{Теория Куммера}

Всюду в этом разделе $n\ge1$~--- целое число, $F$~--- поле, содержащее первообразный корень $\zeta$ степени $n$ из $1$, и
$\Om$~--- фиксированное алгебраическое замыкание $F$; все рассматриваемые расширения $F$ лежат в $\Om$. Мы опишем все
конечные абелевы расширения $F$, экспонента группы Галуа которых делит $n$. Оказывается, все они получаются присоединением
корней степени $n$ из элементов $F$ и взаимно однозначно соответствуют подгруппам $B\subset F^\times$, содержащим группу
$F^{\times n}$ всех $n$-х степеней с конечным индексом. Циклические расширения из \rf{5.27}--\rf{5.28}~--- частный случай: там
присоединяется один корень.

\paragraph{Корни из единицы.} Обозначим через $\mu_n$ группу корней степени $n$ из $1$ в $\Om$. Условие на $F$ равносильно
двум условиям: $\operatorname{char}F\nmid n$ и $\mu_n\subset F$. Действительно, в характеристике $p\mid n$ первообразных
корней степени $n$ из $1$ нет (начало раздела о круговых расширениях), а при $\operatorname{char}F\nmid n$ группа $\mu_n$
циклическая порядка $n$ и её образующие~--- это в точности первообразные корни степени $n$ (\rf{5.7b}, \rf{5.7c}, \rf{5.8}); поэтому
$\zeta\in F$ равносильно $\mu_n=\langle\zeta\rangle\subset F$. Итак, всюду ниже
\[
\mu_n=\{1,\zeta,\dots,\zeta^{n-1}\}\subset F^\times
\]
~--- циклическая группа порядка $n$; для каждого $d\mid n$ в ней ровно одна подгруппа порядка $d$, а именно
$\mu_d=\langle\zeta^{n/d}\rangle$ (\rf{5.7b}(b)), так что $F$ содержит и первообразный корень $\zeta^{n/d}$ степени $d$ из
$1$. Что происходит, если $\mu_n\not\subset F$ или $\operatorname{char}F\mid n$, обсуждается в \rf{5.31i} и \rf{5.31j}.

Главное следствие условия $\mu_n\subset F$: \emph{все корни степени $n$ из одного элемента $a\in F^\times$ отличаются друг
от друга множителями из $F$} (ср.\ \rf{2.10}). Если $\alpha^n=\beta^n=a$, то $(\alpha/\beta)^n=1$, т.\,е.\ $\alpha=\xi\beta$ с
$\xi\in\mu_n\subset F$. Поэтому поле, содержащее $F$ и один корень из $a$, содержит их все, а автоморфизмы над $F$
действуют на $\mu_n$ тождественно.

\paragraph{Экспонента.} \emph{Экспонентой} конечной группы $G$ называется наименьшее $m\ge1$, для которого $g^m=1$ при
всех $g\in G$; это наименьшее общее кратное порядков элементов $G$. Экспонента $G$ делит $n$ тогда и только тогда, когда
$g^n=1$ для всех $g\in G$. Например, $(\Z/2)^r$ имеет экспоненту $2$, а $\Z/4\times\Z/2$~--- экспоненту $4$. Говоря об
экспоненте абелева расширения (\rf{3.16}), мы имеем в виду экспоненту его группы Галуа.

\paragraph{Обозначения.} Положим $F^{\times n}=\{c^n\mid c\in F^\times\}$. Для подгруппы $B$ с $F^{\times n}\subset
B\subset F^\times$ через $F[B^{1/n}]$ обозначим подполе $\Om$, порождённое над $F$ всеми $\alpha\in\Om$ с $\alpha^n\in B$.
По сказанному выше для этого достаточно взять по одному корню $b^{1/n}$ для каждого $b\in B$; здесь и далее $a^{1/n}$
означает \emph{какой-нибудь} корень степени $n$ из $a$ в $\Om$. В частности, поле $F[a^{1/n}]=F[\alpha]$, $\alpha^n=a$,
не зависит от выбора корня (ср.\ абзац перед \rf{5.28}). Для конечного расширения $E\subset\Om$ поля $F$ положим
\[
B(E)=F^\times\cap E^{\times n}=\{a\in F^\times\mid a\ \text{является $n$-й степенью элемента $E$}\}.
\]
Это подгруппа $F^\times$, содержащая $F^{\times n}$. Если $a\in B(E)$, то \emph{все} корни степени $n$ из $a$ лежат в $E$.

\begin{theorem}[5.30: теорема Куммера]
Отображения
\[
E\mapsto B(E)=F^\times\cap E^{\times n},\qquad B\mapsto F[B^{1/n}]
\]
~--- взаимно обратные биекции между множествами
\begin{enumerate}[label=\textup{(\alph*)}]
\item конечных абелевых расширений $E\subset\Om$ поля $F$, экспонента которых делит $n$, и
\item подгрупп $B\subset F^\times$, содержащих $F^{\times n}$ в качестве подгруппы конечного индекса \textup(т.\,е.\
конечных подгрупп $B/F^{\times n}$ группы $F^\times/F^{\times n}$\textup).
\end{enumerate}
Если $E\leftrightarrow B$, то $[E:F]=(B:F^{\times n})$. Обе биекции сохраняют включения.
\end{theorem}

\begin{keepfig}
\[
\begin{tikzcd}[row sep=3.2em]
\Bigl\{\text{конечные абелевы расширения } E\subset\Om \text{ поля } F,\ \text{экспонента которых делит } n\Bigr\}
\arrow[d,cnew,shift left=6,"{\textstyle E\ \mapsto\ B(E)=F^\times\cap E^{\times n}}"]\\
\Bigl\{\text{подгруппы } B,\ F^{\times n}\subset B\subset F^\times,\ (B:F^{\times n})<\infty\Bigr\}
\arrow[u,cold,shift left=6,"{\textstyle B\ \mapsto\ F[B^{1/n}]}"]
\end{tikzcd}
\]
\figcap{Соответствие Куммера. Степень расширения равна индексу подгруппы: $[E:F]=(B:F^{\times n})$.}
\end{keepfig}

\paragraph{План доказательства.} Доказательство опирается на три факта.
\begin{enumerate}[label=(\arabic*),leftmargin=*]
\item Для \emph{любого} конечного расширения Галуа $E/F$ с группой $G$ есть канонический изоморфизм
$B(E)/F^{\times n}\cong\Hom(G,\mu_n)$ (\rf{5.30d}); его сюръективность~--- это теорема Гильберта~90 (\rf{5.23}).
\item Если $G$ абелева и её экспонента делит $n$, то $|\Hom(G,\mu_n)|=|G|$ (\rf{5.30a}). Именно здесь нужно условие на
экспоненту.
\item Поле $F[B^{1/n}]$~--- конечное абелево расширение, экспонента которого делит $n$ (\rf{5.30f}).
\end{enumerate}
После этого теорема получается сравнением порядков. Начнём с нужных фактов о конечных абелевых группах.

\subsection*{Конечные абелевы группы экспоненты, делящей $n$}

\begin{lemma}[5.30a]
Пусть $A$~--- конечная абелева группа, экспонента которой делит $n$.
\begin{enumerate}[label=\textup{(\alph*)}]
\item $A\cong\Z/d_1\times\dots\times\Z/d_r$, где все $d_i$ делят $n$; в частности, $A$ изоморфна подгруппе группы
$(\Z/n)^r$.
\item Группа характеров $\Hom(A,\mu_n)$ \textup(с поточечным умножением\textup) изоморфна $A$; в частности,
$|\Hom(A,\mu_n)|=|A|$.
\item Для каждого $a\ne1$ из $A$ найдётся $\chi\in\Hom(A,\mu_n)$ с $\chi(a)\ne1$.
\end{enumerate}
\end{lemma}

\begin{proof}
(a) По структурной теореме $A\cong\Z/d_1\times\dots\times\Z/d_r$. Образующая $i$-го сомножителя~--- элемент порядка
$d_i$, а порядок каждого элемента делит экспоненту; значит, $d_i\mid n$. (То же на языке примарного разложения: если
$A\cong\bigoplus\Z/p^{k}$, то в $\Z/p^k$ есть элемент порядка $p^k$, поэтому $p^k\mid n$.) При $d\mid n$ отображение
$\Z/d\to\Z/n$, $x\mapsto x\cdot\frac nd$, ~--- инъективный гомоморфизм: $x\frac nd\equiv0\pmod n\iff d\mid x$.
Произведение таких вложений вкладывает $A$ в $(\Z/n)^r$.

(b) Характер $\chi\colon\Z/d\to\mu_n$ определяется значением $x=\chi(1)$, и $x$ может быть любым элементом $\mu_n$ с
$x^d=1$, т.\,е.\ любым элементом $\mu_d$. Поэтому $\Hom(\Z/d,\mu_n)\cong\mu_d\cong\Z/d$; здесь существенно, что
$d\mid n$: в $\mu_n$ ровно $d$ корней степени $d$ из $1$. Далее, характер произведения~--- это пара характеров
сомножителей: $\Hom(A_1\times A_2,\mu_n)\cong\Hom(A_1,\mu_n)\times\Hom(A_2,\mu_n)$, $\chi\mapsto(\chi|_{A_1},\chi|_{A_2})$,
и обратно $\chi(a_1,a_2)=\chi_1(a_1)\chi_2(a_2)$. Значит, $\Hom(A,\mu_n)\cong\prod_i\Z/d_i\cong A$.

(c) Запишем $A=\Z/d_1\times\dots\times\Z/d_r$ аддитивно и $a=(a_1,\dots,a_r)$; пусть $a_j\ne0$ в $\Z/d_j$. Положим
$\chi(x_1,\dots,x_r)=\zeta_j^{\,x_j}$, где $\zeta_j=\zeta^{n/d_j}$~--- первообразный корень степени $d_j$. Тогда
$\chi(a)=\zeta_j^{\,a_j}\ne1$.
\end{proof}

Изоморфизм в (b) не канонический: он зависит от разложения $A$ в произведение циклических групп и от выбора их
образующих (и $\zeta$). Условие на экспоненту существенно: например, $|\Hom(\Z/4,\mu_2)|=2\ne4$.

\begin{definition}[5.30b]
Пусть $A$ и $G$~--- абелевы группы, в которых $x^n=1$ для всех элементов $x$ (для конечных групп это значит, что
экспонента делит $n$). \emph{Спаривание}~--- это отображение
$\langle\cdot,\cdot\rangle\colon A\times G\to\mu_n$, мультипликативное по каждому аргументу:
\[
\langle aa',\sigma\rangle=\langle a,\sigma\rangle\langle a',\sigma\rangle,\qquad
\langle a,\sigma\tau\rangle=\langle a,\sigma\rangle\langle a,\tau\rangle.
\]
Оно задаёт гомоморфизмы $\lambda\colon A\to\Hom(G,\mu_n)$, $a\mapsto\langle a,\cdot\rangle$, и
$\rho\colon G\to\Hom(A,\mu_n)$, $\sigma\mapsto\langle\cdot,\sigma\rangle$; их ядра называются \emph{левым} и \emph{правым
ядром} спаривания. Спаривание \emph{невырождено}, если оба ядра тривиальны, и \emph{совершенно}, если $\lambda$ и
$\rho$~--- изоморфизмы. \emph{Аннуляторы} подгрупп $A'\subset A$ и $H\subset G$~--- это
\[
A'^{\perp}=\{\sigma\in G\mid\langle a,\sigma\rangle=1\ \text{для всех}\ a\in A'\},\qquad
H^{\perp}=\{a\in A\mid\langle a,\sigma\rangle=1\ \text{для всех}\ \sigma\in H\}.
\]
\end{definition}

Образец~--- спаривание $V\times V^*\to k$ конечномерного векторного пространства с двойственным. Примеры:
$\Z/n\times\Z/n\to\mu_n$, $\langle a,b\rangle=\zeta^{ab}$, и <<скалярное произведение>> $(\Z/n)^r\times(\Z/n)^r\to\mu_n$,
$\langle a,b\rangle=\zeta^{a_1b_1+\dots+a_rb_r}$, совершенны (это видно из следующей леммы: ядра тривиальны). Спаривание
$\Z/4\times\Z/4\to\mu_4$, $\langle a,b\rangle=(-1)^{ab}$, не совершенно: его левое ядро~--- $\{0,2\}$.

\begin{lemma}[5.30c]
Пусть $A$, $G$~--- конечные абелевы группы, экспоненты которых делят $n$, и $\langle\cdot,\cdot\rangle\colon A\times
G\to\mu_n$~--- спаривание.
\begin{enumerate}[label=\textup{(\alph*)}]
\item Если оба ядра тривиальны, то спаривание совершенно и $|A|=|G|$.
\item Пусть спаривание совершенно и $A'\subset A$~--- подгруппа. Тогда формула $(a,\sigma A'^\perp)\mapsto\langle
a,\sigma\rangle$ корректно задаёт совершенное спаривание $A'\times G/A'^\perp\to\mu_n$. В частности, $\lambda$
отображает $A'$ изоморфно \emph{на} подгруппу $\Hom(G/A'^\perp,\mu_n)\subset\Hom(G,\mu_n)$ характеров, тривиальных на
$A'^\perp$, и $|A'^\perp|=|G|/|A'|$. Кроме того, $(A'^\perp)^\perp=A'$, и $A'\mapsto A'^\perp$~--- обращающая
включения биекция между подгруппами $A$ и подгруппами $G$ \textup(обратная биекция~--- $H\mapsto H^\perp$\textup).
\end{enumerate}
\end{lemma}

\begin{proof}
(a) Отображения $\lambda$ и $\rho$ инъективны, поэтому по \rf{5.30a}(b)
\[
|A|\le|\Hom(G,\mu_n)|=|G|\le|\Hom(A,\mu_n)|=|A|.
\]
Значит, всюду равенства, и инъективные отображения между конечными множествами одинакового размера биективны.

(b) Если $\tau\in A'^\perp$ и $a\in A'$, то $\langle a,\sigma\tau\rangle=\langle a,\sigma\rangle\langle a,\tau\rangle=
\langle a,\sigma\rangle$, так что формула корректна; мультипликативность очевидна. Левое ядро нового спаривания
тривиально: если $\langle a,\sigma\rangle=1$ для всех $\sigma\in G$, то $a=1$, так как $\lambda$ инъективно. Правое ядро
тривиально по определению $A'^\perp$. Группа $G/A'^\perp$ имеет экспоненту, делящую $n$, и по (a) новое спаривание
совершенно, $|A'|=|G/A'^\perp|$. Характеры $G/A'^\perp$~--- это характеры $G$, тривиальные на $A'^\perp$, и
изоморфизм $A'\to\Hom(G/A'^\perp,\mu_n)$~--- это ограничение $\lambda$ на $A'$.

Ясно, что $A'\subset(A'^\perp)^\perp$. Спаривание $G\times A\to\mu_n$, $(\sigma,a)\mapsto\langle a,\sigma\rangle$, тоже
совершенно, и по уже доказанному для подгруппы $H=A'^\perp\subset G$ имеем $|H^\perp|=|A|/|H|=|A|\,|A'|/|G|=|A'|$ (так как
$|A|=|G|$). Значит, $(A'^\perp)^\perp=A'$; так же $(H^\perp)^\perp=H$ для $H\subset G$. Отсюда биективность;
включения, очевидно, обращаются.
\end{proof}

Это та же двойственность, что в линейной алгебре: аннулятор аннулятора подпространства $W\subset V^*$ равен $W$. Роль
размерностей играют порядки, а роль равенства $\dim V^*=\dim V$~--- лемма \rf{5.30a}(b).

\subsection*{Изоморфизм Куммера}

\begin{proposition}[5.30d]
Пусть $E/F$~--- конечное расширение Галуа с группой $G$ \textup(не обязательно абелевой\textup).
\begin{enumerate}[label=\textup{(\alph*)}]
\item Пусть $a\in B(E)$ и $\alpha\in E$, $\alpha^n=a$. Отображение $\chi_a\colon G\to\mu_n$, $\sigma\mapsto
\sigma(\alpha)/\alpha$, корректно определено, не зависит от выбора корня $\alpha$ и является гомоморфизмом групп.
\item Отображение $a\mapsto\chi_a$~--- сюръективный гомоморфизм $B(E)\to\Hom(G,\mu_n)$ с ядром $F^{\times n}$.
Следовательно, он индуцирует изоморфизм
\[
\kappa_E\colon B(E)/F^{\times n}\xrightarrow{\ \sim\ }\Hom(G,\mu_n),\qquad
a\mapsto\Bigl(\sigma\mapsto\frac{\sigma(a^{1/n})}{a^{1/n}}\Bigr).
\]
\end{enumerate}
\end{proposition}

\begin{proof}
(a) $(\sigma\alpha)^n=\sigma(\alpha^n)=\sigma(a)=a=\alpha^n$, поэтому $\sigma(\alpha)/\alpha\in\mu_n$. Любой другой корень
из $a$ имеет вид $\xi\alpha$ с $\xi\in\mu_n\subset F$, и $\sigma(\xi\alpha)/(\xi\alpha)=\xi\sigma(\alpha)/(\xi\alpha)=
\sigma(\alpha)/\alpha$, так как $\sigma$ тождественно на $F$. Гомоморфность: $\tau\alpha$~--- тоже корень из $a$, поэтому
по уже доказанному $\sigma(\tau\alpha)/\tau\alpha=\chi_a(\sigma)$, и
\[
\chi_a(\sigma\tau)=\frac{\sigma(\tau\alpha)}{\alpha}=\frac{\sigma(\tau\alpha)}{\tau\alpha}\cdot\frac{\tau\alpha}{\alpha}
=\chi_a(\sigma)\,\chi_a(\tau).
\]
Условие $\mu_n\subset F$ использовано здесь дважды: при проверке независимости от выбора корня и (что то же самое) в том,
что $G$ действует на значениях $\chi_a$ тривиально.

(b) Если $\alpha^n=a$ и $\beta^n=b$, то $(\alpha\beta)^n=ab$ и $\sigma(\alpha\beta)/(\alpha\beta)=\chi_a(\sigma)\chi_b(\sigma)$;
значит, $\chi_{ab}=\chi_a\chi_b$.

\emph{Ядро.} $\chi_a=1$ $\iff$ $\sigma(\alpha)=\alpha$ для всех $\sigma\in G$ $\iff$ $\alpha\in E^G=F$. Если $\alpha\in F$, то
$a=\alpha^n\in F^{\times n}$; обратно, если $a=c^n$ с $c\in F^\times$, то в качестве $\alpha$ можно взять $c$, и
$\chi_a=1$.

\emph{Сюръективность.} Пусть $\chi\colon G\to\mu_n$~--- гомоморфизм. Так как $G$ действует на $\mu_n\subset F$
тривиально, $\chi(\sigma\tau)=\chi(\sigma)\chi(\tau)=\chi(\sigma)\cdot\sigma\bigl(\chi(\tau)\bigr)$, т.\,е.\ $\chi$~---
скрещённый гомоморфизм $G\to E^\times$. По теореме Гильберта~90 (\rf{5.23}) он главный: найдётся $\beta\in E^\times$ с
$\chi(\sigma)=\sigma(\beta)/\beta$ для всех $\sigma\in G$. Положим $a=\beta^n$. Тогда
$\sigma(a)=\sigma(\beta)^n=\chi(\sigma)^n\beta^n=a$ для всех $\sigma$, т.\,е.\ $a\in E^G=F$. Итак, $a\in F^\times\cap
E^{\times n}=B(E)$, и, взяв $\alpha=\beta$, получаем $\chi_a=\chi$.
\end{proof}

\paragraph{Когомологический смысл.} Предложение \rf{5.30d}~--- кусок длинной точной последовательности когомологий. Возведение
в степень $n$ даёт короткую точную последовательность $G$-модулей (\rf{5.20a}(b); здесь $\mu_n(E)=\mu_n$, так как
$\mu_n\subset F\subset E$)
\[
1\longrightarrow\mu_n\longrightarrow E^\times\xrightarrow{\ x\mapsto x^n\ }E^{\times n}\longrightarrow1,
\]
и точная последовательность \rf{5.21b} для неё даёт верхнюю строку диаграммы
\par\medskip
\begin{keepfig}
\[
\begin{tikzcd}[column sep=1.8em,row sep=2.2em]
1\arrow[r] & \mu_n^G\arrow[r]\arrow[d,equal] & (E^\times)^G\arrow[r,"x\mapsto x^n"]\arrow[d,equal] & (E^{\times n})^G\arrow[r,"\delta"]
\arrow[d,equal] & H^1(G,\mu_n)\arrow[r]\arrow[d,equal] & H^1(G,E^\times)\arrow[d,equal,"\text{(5.23)}"]\\
1\arrow[r] & \mu_n\arrow[r] & F^\times\arrow[r,"x\mapsto x^n"] & F^\times\cap E^{\times n}\arrow[r,cnew,"\delta"]
& \Hom(G,\mu_n)\arrow[r] & 1
\end{tikzcd}
\]
\figcap{Длинная точная последовательность для $1\to\mu_n\to E^\times\to E^{\times n}\to1$ и её члены.}
\end{keepfig}

Здесь $\mu_n^G=\mu_n$, так как $G$ действует на $\mu_n\subset F$ тривиально; $(E^\times)^G=F^\times$ (\rf{5.20}); $(E^{\times n})^G=
F^\times\cap E^{\times n}$, так как $n$-я степень $x^n$ неподвижна относительно $G$ ровно тогда, когда $x^n\in F$;
$H^1(G,E^\times)=1$ по теореме Гильберта~90 (\rf{5.23}); наконец, $H^1(G,\mu_n)=\Hom(G,\mu_n)$, так как при тривиальном
действии скрещённые гомоморфизмы~--- это обычные гомоморфизмы, а главные скрещённые гомоморфизмы
$\sigma\mapsto\sigma(\xi)/\xi$ тривиальны (\rf{5.21}(c)). Получаем точную последовательность
\[
1\to\mu_n\to F^\times\xrightarrow{\ x\mapsto x^n\ }F^\times\cap E^{\times n}\xrightarrow{\ \delta\ }\Hom(G,\mu_n)\to1,
\]
т.\,е.\ $\Hom(G,\mu_n)\cong H^1(G,\mu_n)\cong\operatorname{coker}\bigl(F^\times\xrightarrow{x\mapsto x^n}F^\times\cap
E^{\times n}\bigr)=B(E)/F^{\times n}$. Связывающий гомоморфизм $\delta$ устроен так (\rf{5.21b}): для $a\in(E^{\times n})^G$
берём прообраз $y\in E^\times$, $y^n=a$; тогда $\sigma\mapsto\sigma(y)/y$~--- скрещённый гомоморфизм со значениями в
$\mu_n$, и $\delta(a)$~--- его класс. Это в точности $\chi_a$ (ср.\ \rf{5.27b}), а доказательство \rf{5.30d}~--- проверка точности
этой последовательности <<вручную>>. В частности, сюръективность $\kappa_E$~--- это равенство $H^1(G,E^\times)=1$.

\begin{aside}[5.30e: скрещённые гомоморфизмы и гомоморфизмы]
Чем скрещённые гомоморфизмы отличаются от гомоморфизмов, разобрано в \rf{5.21a}. В теории Куммера значения $\chi_a$ лежат в
$\mu_n\subset F$, где $G$ действует тривиально, поэтому $H^1(G,\mu_n)=\Hom(G,\mu_n)$ и появляются обычные гомоморфизмы
(\rf{5.21}(c)); в теореме Гильберта~90 значения лежат в $E^\times$, и скрещённость существенна.
\end{aside}

\subsection*{Поле $F[B^{1/n}]$}

\begin{lemma}[5.30f]
Пусть $F^{\times n}\subset B\subset F^\times$, $(B:F^{\times n})<\infty$, классы элементов $a_1,\dots,a_k\in B$
порождают $B/F^{\times n}$, и $\alpha_i\in\Om$, $\alpha_i^n=a_i$. Положим $E_B=F[B^{1/n}]$.
\begin{enumerate}[label=\textup{(\alph*)}]
\item $E_B=F[\alpha_1,\dots,\alpha_k]$.
\item Каждое $F[\alpha_i]/F$~--- расширение Галуа, и $\sigma\mapsto\sigma(\alpha_i)/\alpha_i$~--- инъективный гомоморфизм
$\Gal(F[\alpha_i]/F)\to\mu_n$; в частности, $F[\alpha_i]/F$ циклично и его степень делит $n$.
\item $E_B/F$~--- конечное расширение Галуа, и $\sigma\mapsto\bigl(\sigma(\alpha_1)/\alpha_1,\dots,
\sigma(\alpha_k)/\alpha_k\bigr)$~--- инъективный гомоморфизм $\Gal(E_B/F)\to\mu_n^k$. Следовательно, $E_B/F$~---
конечное абелево расширение, экспонента которого делит $n$.
\end{enumerate}
\end{lemma}

\begin{proof}
(a) Каждое $b\in B$ имеет вид $b=a_1^{e_1}\cdots a_k^{e_k}c^n$ с $e_i\in\Z$, $c\in F^\times$. Тогда
$\alpha_1^{e_1}\cdots\alpha_k^{e_k}c$~--- корень степени $n$ из $b$, а остальные корни отличаются от него множителями из
$\mu_n\subset F$; все они лежат в $F[\alpha_1,\dots,\alpha_k]$. Обратное включение очевидно.

(b) Корни многочлена $X^n-a_i$~--- это $n$ различных элементов $\xi\alpha_i$, $\xi\in\mu_n$ (здесь $a_i\ne0$ и
$|\mu_n|=n$), и все они лежат в $F[\alpha_i]$. Значит, $F[\alpha_i]$~--- поле разложения сепарабельного многочлена, т.\,е.\
расширение Галуа (\rf{3.10}). Отображение $\sigma\mapsto\sigma(\alpha_i)/\alpha_i$~--- это гомоморфизм $\chi_{a_i}$ из \rf{5.30d}(a)
для $E=F[\alpha_i]$; если $\sigma(\alpha_i)=\alpha_i$, то $\sigma$ тождественно на $F[\alpha_i]$. Подгруппы циклической
группы $\mu_n$ цикличны, и их порядки делят $n$.

(c) $E_B$~--- композит расширений Галуа $F[\alpha_i]$, поэтому оно конечно и Галуа над $F$ (\rf{3.21}). Координаты
отображения~--- гомоморфизмы $\chi_{a_i}$ для $E=E_B$ (\rf{5.30d}(a)). Если $\sigma(\alpha_i)=\alpha_i$ для всех $i$, то
$\sigma$ тождественно на $F[\alpha_1,\dots,\alpha_k]=E_B$. Группа $\mu_n^k$ абелева, и $x^n=1$ для всех её элементов; эти
свойства наследует любая подгруппа.
\end{proof}

\paragraph{Расслоённые произведения.} Отображение из (c) раскладывается в композицию: сначала вложение
$\Gal(E_B/F)\hookrightarrow\prod_i\Gal(F[\alpha_i]/F)$, $\sigma\mapsto(\sigma|_{F[\alpha_i]})_i$ (ограничения определены,
так как $F[\alpha_i]$ нормальны над $F$), затем вложения $\Gal(F[\alpha_i]/F)\hookrightarrow\mu_n$ из (b). Образ,
вообще говоря,~--- собственная подгруппа произведения. Для двух сомножителей он описан в \rf{3.21}: группа Галуа композита $E_1E_2$ двух расширений Галуа~---
расслоённое произведение
\[
\Gal(E_1E_2/F)\cong\Gal(E_1/F)\times_{\Gal(E_1\cap E_2/F)}\Gal(E_2/F)=\{(\sigma_1,\sigma_2)\mid
\sigma_1|_{E_1\cap E_2}=\sigma_2|_{E_1\cap E_2}\}.
\]
Например, при $F=\Q$, $n=2$ и $a_1=2$, $a_2=3$, $a_3=6$ образ $\Gal(\Q[\sqrt2,\sqrt3]/\Q)$ в $\mu_2^3$~--- это подгруппа
$\{(s,t,st)\}$ порядка $4$, а не всё $\mu_2^3$. Для вывода (c) точное описание образа не нужно: достаточно того, что
это подгруппа абелевой группы, экспонента которой делит $n$.

\subsection*{Доказательство теоремы \rf{5.30}}

\emph{Шаг 1: отображения определены, и $[E:F]=(B(E):F^{\times n})$.} Пусть $E$ из множества (a), $G=\Gal(E/F)$. По \rf{5.30d}
и \rf{5.30a}(b)
\[
(B(E):F^{\times n})=|\Hom(G,\mu_n)|=|G|=[E:F]<\infty.\tag{$*$}
\]
Значит, $B(E)$ принадлежит множеству (b). Обратно, если $B$ из множества (b), то $F[B^{1/n}]$ принадлежит множеству (a)
по \rf{5.30f}(c).

\emph{Шаг 2: $F[B(E)^{1/n}]=E$ для $E$ из (a).} Положим $E'=F[B(E)^{1/n}]$. Если $\alpha\in\Om$ и $\alpha^n=a\in B(E)$, то
$a=y^n$ с $y\in E$ и $\alpha=\xi y$ с $\xi\in\mu_n\subset F$; значит, $\alpha\in E$, и $E'\subset E$. Каждый элемент $B(E)$
имеет корень степени $n$ в $E'$, т.\,е.\ $B(E)\subset B(E')$. По шагу 1 $E'$ принадлежит (a), и формула $(*)$ применима к
$E$ и к $E'$:
\[
[E:F]\ \ge\ [E':F]\ \overset{(*)}{=}\ (B(E'):F^{\times n})\ \ge\ (B(E):F^{\times n})\ \overset{(*)}{=}\ [E:F].
\]
Значит, всюду равенства, $[E':F]=[E:F]$, и из $E'\subset E$ следует $E'=E$.

\medskip
\begin{keepfig}
\[
\begin{tikzcd}[column sep=4em,row sep=2.6em]
E\arrow[r,cold,mapsto] & B(E)\arrow[d,hook]\\
E'=F[B(E)^{1/n}]\arrow[u,hook]\arrow[r,cnew,mapsto] & B(E')
\end{tikzcd}
\]
\figcap{Шаг 2: $E'\subset E$, но $B(E)\subset B(E')$. Формула $(*)$ превращает эти включения в неравенства, которые
замыкаются в круг.}
\end{keepfig}

\emph{Шаг 3: $B(F[B^{1/n}])=B$ для $B$ из (b).} Положим $E=F[B^{1/n}]$, $G=\Gal(E/F)$. Ясно, что $B\subset B(E)$.
Для $\sigma\in G$ отображение $b\mapsto\chi_b(\sigma)$ (обозначения \rf{5.30d} для $E$)~--- гомоморфизм $B\to\mu_n$,
тривиальный на $F^{\times n}$ (\rf{5.30d}(b)); получаем гомоморфизм
\[
G\to\Hom(B/F^{\times n},\mu_n),\qquad\sigma\mapsto\bigl(b\mapsto\chi_b(\sigma)\bigr).
\]
Он инъективен: если $\chi_b(\sigma)=1$ для всех $b\in B$, то $\sigma$ оставляет на месте все корни $b^{1/n}$, $b\in B$, а
они порождают $E$; значит, $\sigma=1$. Группа $B/F^{\times n}$ конечна, и её экспонента делит $n$ ($b^n\in
F^{\times n}$). По \rf{5.30a}(b) и $(*)$
\[
[E:F]=|G|\ \le\ |\Hom(B/F^{\times n},\mu_n)|=(B:F^{\times n})\ \le\ (B(E):F^{\times n})\overset{(*)}{=}[E:F].
\]
Значит, $(B:F^{\times n})=(B(E):F^{\times n})$, и из $B\subset B(E)$ следует $B=B(E)$.

\emph{Шаг 4.} По шагам 2 и 3 отображения взаимно обратны. Формула для степени~--- это $(*)$. Включения сохраняются:
если $E_1\subset E_2$, то $B(E_1)\subset B(E_2)$, а если $B_1\subset B_2$, то $F[B_1^{1/n}]\subset F[B_2^{1/n}]$.
\qed

\begin{remark}[5.30g: конец доказательства на языке спариваний]
Шаг 3~--- утверждение о спаривании $B/F^{\times n}\times G\to\mu_n$, $\langle b,\sigma\rangle=\chi_b(\sigma)$. Его левое
ядро тривиально (по \rf{5.30d} отображение $\kappa_E$ инъективно на $B(E)/F^{\times n}\supset B/F^{\times n}$), правое ядро
тривиально (автоморфизм, не двигающий все $b^{1/n}$, тождествен на $E$). По \rf{5.30c}(a) оно совершенно, откуда
$|B/F^{\times n}|=|G|=|B(E)/F^{\times n}|$.

Можно рассуждать и так. Пусть $X=\kappa_E(B/F^{\times n})\subset\Hom(G,\mu_n)$~--- множество характеров $\chi_b$, $b\in B$,
и $H=\bigcap_{b\in B}\ker\chi_b$~--- их общее ядро. Каждый $\chi\in X$ тривиален на $H$, так что $X\subset
\Hom(G/H,\mu_n)$. \emph{На самом деле $\kappa_E$ отображает $B/F^{\times n}$ изоморфно на $\Hom(G/H,\mu_n)$.}
Действительно, спаривание $X\times G/H\to\mu_n$, $(\chi,\sigma H)\mapsto\chi(\sigma)$, корректно и невырождено: характер из
$X$, равный $1$ на всей $G$, тривиален, а $\sigma$ с $\chi(\sigma)=1$ для всех $\chi\in X$ лежит в $H$ по определению. По
\rf{5.30c}(a) $|X|=|G/H|=|\Hom(G/H,\mu_n)|$, и включение $X\subset\Hom(G/H,\mu_n)$~--- равенство. Это и есть утверждение о
двойном аннуляторе из \rf{5.30c}(b). Остаётся заметить, что $H=1$: элемент $H$ оставляет на месте все $b^{1/n}$, $b\in B$,
и потому тождествен на $E$. Значит, $\kappa_E(B/F^{\times n})=\Hom(G,\mu_n)=\kappa_E(B(E)/F^{\times n})$, и, так как
$\kappa_E$ инъективно, $B=B(E)$.
\end{remark}

\begin{remark}[5.30h: экспонента ровно $n$]
Некоторым $B$ соответствуют расширения, экспонента которых меньше $n$: подгруппе $B=F^{\times n}$ соответствует $E=F$, а
при $F=\Q[i]$, $n=4$ подгруппе $B=\langle9\rangle F^{\times4}$ соответствует $E=\Q[i,\sqrt3]$ с группой $\Z/2$
(\rf{5.31d}). Поэтому в (a) берутся расширения, экспонента которых \emph{делит} $n$. Если нужна экспонента ровно $n$, то
теорема ограничивается на биекцию между абелевыми расширениями экспоненты ровно $n$ и подгруппами $B$, для которых
экспонента $B/F^{\times n}$ равна $n$: группы $B(E)/F^{\times n}\cong\Hom(G,\mu_n)$ и $G$ изоморфны (\rf{5.30d}, \rf{5.30a}(b)) и
потому имеют одинаковую экспоненту.
\end{remark}

\subsection*{Следствия}

\begin{corollary}[5.30i: двучленные уравнения]
Пусть $a\in F^\times$, $\alpha^n=a$, и $d$~--- порядок класса $a$ в $F^\times/F^{\times n}$.
\begin{enumerate}[label=\textup{(\alph*)}]
\item $d\mid n$; $\alpha^d\in F$, и никакая меньшая положительная степень $\alpha$ не лежит в $F$; $F[\alpha]/F$~---
циклическое расширение степени $d$, а $X^d-\alpha^d$~--- минимальный многочлен $\alpha$ над $F$.
\item $X^n-a$ неприводим над $F$ $\iff$ класс $a$ имеет порядок $n$ $\iff$ $a$ не является $p$-й степенью в $F$ ни для
какого простого $p\mid n$.
\item Если $n=\ell$ простое, то $X^\ell-a$ либо неприводим, либо раскладывается в $F[X]$ на линейные множители.
\end{enumerate}
\end{corollary}

\begin{proof}
(a) Множество $\{m\in\Z\mid\alpha^m\in F\}$~--- подгруппа $\Z$ (так как $\alpha\ne0$), и $\alpha^m\in F\iff a^m\in
F^{\times n}$: если $\alpha^m=c\in F$, то $a^m=c^n$; если $a^m=c^n$ с $c\in F^\times$, то $(\alpha^m/c)^n=1$ и
$\alpha^m\in c\mu_n\subset F$. Значит, эта подгруппа равна $d\Z$. Она содержит $n$, так что $d\mid n$. (Иначе говоря:
если $\alpha^d\in F$ и $n=kd+r$, $0\le r<d$, то $\alpha^r=\alpha^n(\alpha^d)^{-k}\in F$, и минимальность $d$ даёт $r=0$.)
Поле $F$ содержит первообразный корень $\zeta^{n/d}$ степени $d$ из $1$, поэтому по \rf{5.27}, применённому к $d$ вместо $n$,
$F[\alpha]/F$ циклично степени $d$. Многочлен $X^d-\alpha^d\in F[X]$ степени $d=[F[\alpha]:F]$ имеет корень $\alpha$,
значит, он минимальный. (То же даёт теорема \rf{5.30} для $B=\langle a\rangle F^{\times n}$: $(B:F^{\times n})=d$ и
$F[B^{1/n}]=F[\alpha]$ по \rf{5.30f}(a).)

(b) Первая равносильность: $X^n-a$ неприводим $\iff$ $[F[\alpha]:F]=n$ $\iff$ $d=n$ (по (a)). Равносильность
неприводимости $X^n-a$ и последнего условия~--- это \rf{5.29}(a); вместе с первой равносильностью она даёт вторую.

(c) Здесь $d\in\{1,\ell\}$. При $d=\ell$ многочлен неприводим; при $d=1$ имеем $\alpha\in F$, и все корни $\xi\alpha$,
$\xi\in\mu_\ell$, лежат в $F$.
\end{proof}

Условие <<$a$ не является $n$-й степенью>> для неприводимости $X^n-a$ недостаточно: над $\Q[i]$ число $4$ не является
четвёртой степенью, но $X^4-4=(X^2-2)(X^2+2)$ (\rf{5.31d}; ср.\ \rf{5.29}(a)). Без корней из единицы утверждение (c) неверно: над $\Q$ многочлен
$X^3-8=(X-2)(X^2+2X+4)$ имеет ровно один корень.

\begin{corollary}[5.30j: какие корни лежат в радикальном расширении]
Пусть $a_1,\dots,a_k\in F^\times$, $E=F[a_1^{1/n},\dots,a_k^{1/n}]$ и $B=\langle a_1,\dots,a_k\rangle F^{\times n}$.
\begin{enumerate}[label=\textup{(\alph*)}]
\item $[E:F]$ равно порядку подгруппы группы $F^\times/F^{\times n}$, порождённой классами $a_1,\dots,a_k$.
\item Для $b\in F^\times$: $b^{1/n}\in E$ $\iff$ $b\in B$, т.\,е.\ $b=a_1^{e_1}\cdots a_k^{e_k}c^n$ с $e_i\in\Z$,
$c\in F^\times$.
\item $F[a^{1/n}]=F[b^{1/n}]$ $\iff$ классы $a$ и $b$ порождают одну и ту же подгруппу $F^\times/F^{\times n}$, т.\,е.\
$b=a^rc^n$, где $c\in F^\times$, а $r$ взаимно просто с порядком класса $a$.
\end{enumerate}
\end{corollary}

\begin{proof}
По \rf{5.30f}(a) $E=F[B^{1/n}]$, и по теореме \rf{5.30} $B(E)=B$ и $[E:F]=(B:F^{\times n})$; это (a). (b): $b^{1/n}\in E\iff b\in
B(E)=B$. (c): $F[a^{1/n}]=F[(\langle a\rangle F^{\times n})^{1/n}]$, и по биективности соответствия эти поля совпадают ровно
тогда, когда совпадают подгруппы $\langle a\rangle F^{\times n}$ и $\langle b\rangle F^{\times n}$. В циклической группе
порядка $d$ элемент $\bar a^r$ порождает ту же подгруппу, что $\bar a$, тогда и только тогда, когда $(r,d)=1$.

Направление <<$\Leftarrow$>> в (c) проверяется и непосредственно, как <<если>> в \rf{5.28}, только вместо $n$ берётся
порядок $d$ класса $a$ и используется $\alpha^d\in F$ (\rf{5.30i}(a)).
\end{proof}

При $\operatorname{ord}\bar a=n$ это предложение \rf{5.28}; теперь степень может быть любой.

\begin{corollary}[5.30k: промежуточные поля]
Пусть $B$ из множества \textup{(b)} теоремы \rf{5.30} и $E=F[B^{1/n}]$. Промежуточные поля $F\subset K\subset E$ взаимно
однозначно соответствуют подгруппам $B'$ с $F^{\times n}\subset B'\subset B$: $K\mapsto B(K)$, $B'\mapsto F[B'^{1/n}]$,
причём $[K:F]=(B':F^{\times n})$. В частности, промежуточных полей степени $d$ столько же, сколько подгрупп порядка $d$ в
$B/F^{\times n}$.
\end{corollary}

\begin{proof}
Группа $G=\Gal(E/F)$ абелева, поэтому каждая её подгруппа нормальна, и по основной теореме (\rf{3.17}) каждое $K$ является
расширением Галуа поля $F$ с группой $G/\Gal(E/K)$~--- абелевой, экспонента которой делит $n$. Значит, $K$ лежит в
множестве (a), и $B(K)\subset B(E)=B$. Обратно, $F[B'^{1/n}]\subset F[B^{1/n}]=E$. Остаётся ограничить биекцию \rf{5.30},
сохраняющую включения.
\end{proof}

\begin{corollary}[5.30l: циклические расширения]
Пусть $d\mid n$. Циклические расширения $E\subset\Om$ поля $F$ степени $d$ взаимно однозначно соответствуют циклическим
подгруппам порядка $d$ в $F^\times/F^{\times n}$: подгруппе $\langle\bar a\rangle$ соответствует $F[a^{1/n}]$. В
частности, каждое циклическое расширение $E/F$ степени $d$ имеет вид $E=F[\alpha]$, где $\alpha^d\in F$.
\end{corollary}

\begin{proof}
Циклическое расширение степени $d$ абелево, и его экспонента $d$ делит $n$. По \rf{5.30d} и \rf{5.30a}(b) группа
$B(E)/F^{\times n}\cong\Hom(\Z/d,\mu_n)\cong\Z/d$ циклична порядка $d$; пусть её порождает класс $a$. Тогда $\langle a\rangle F^{\times n}=B(E)$, и по \rf{5.30f}(a) и
теореме \rf{5.30} $F[a^{1/n}]=F[B(E)^{1/n}]=E$; для $\alpha=a^{1/n}$ имеем $\alpha^d\in F$ по \rf{5.30i}(a). Обратно, по \rf{5.30i}(a)
циклической подгруппе порядка $d$ соответствует циклическое расширение степени $d$. Соответствие биективно как ограничение
биекции \rf{5.30}: группа $B/F^{\times n}\cong\Hom(\Gal(E/F),\mu_n)$ циклична порядка $d$ тогда и только тогда, когда
$\Gal(E/F)$ циклична порядка $d$ (\rf{5.30a}(b)).
\end{proof}

При $d=n$ это \rf{5.28a}: циклические расширения степени $n$ соответствуют циклическим подгруппам порядка $n$ в
$F^\times/F^{\times n}$. Нужный элемент $\alpha$ строится явно, как в доказательстве \rf{5.27} и абзаце после него, только с $d$
вместо $n$ и первообразным корнем $\zeta_d=\zeta^{n/d}$ вместо $\zeta$: если $\sigma$ порождает $\Gal(E/F)$, то годится
резольвента Лагранжа $\alpha=\sum_{i=0}^{d-1}\zeta_d^{\,i}\sigma^i(\gamma)$ для любого $\gamma\in E$, при котором она
отлична от нуля (такие $\gamma$ есть по \rf{5.14}, \rf{5.15a}); для неё $\sigma\alpha=\zeta_d^{-1}\alpha$.

\begin{corollary}[5.30m: наибольшее подрасширение экспоненты, делящей $n$]
Пусть $E/F$~--- любое конечное расширение Галуа с группой $G$. Тогда $E_0=F[B(E)^{1/n}]$~--- наибольшее подполе $E$,
абелево над $F$ с экспонентой, делящей $n$, и $[E_0:F]=(B(E):F^{\times n})=|\Hom(G,\mu_n)|$. Если $N\subset G$~---
подгруппа, порождённая коммутаторами и $n$-ми степенями, то $E_0=E^N$ и $[E_0:F]=(G:N)$.
\end{corollary}

\begin{proof}
По \rf{5.30d} индекс $(B(E):F^{\times n})=|\Hom(G,\mu_n)|$ конечен, поэтому $E_0$ лежит в множестве (a) (\rf{5.30f}), и
$E_0\subset E$ (как в шаге 2). Если $K\subset E$ абелево над $F$ с экспонентой, делящей $n$, то по теореме \rf{5.30}
$K=F[B(K)^{1/n}]\subset F[B(E)^{1/n}]=E_0$, так как $B(K)\subset B(E)$. По теореме $B(E_0)=B(E)$, откуда
$[E_0:F]=(B(E):F^{\times n})$. Далее, $N$ нормальна (порождающее множество инвариантно относительно сопряжений), $G/N$
абелева и её экспонента делит $n$; каждый характер $G\to\mu_n$ тривиален на $N$, поэтому $|\Hom(G,\mu_n)|=
|\Hom(G/N,\mu_n)|=(G:N)$. Поле $E^N$ абелево над $F$ с группой $G/N$, так что $E^N\subset E_0$, а степени равны.
\end{proof}

\subsection*{Примеры}

\begin{example}[5.31]
(a) $F=\R$, $n=2$. Каждое положительное вещественное число~--- квадрат, поэтому $\R^\times/\R^{\times2}=\{\pm1\}$. Его
подгруппам соответствуют $\R$ и $\R[\sqrt{-1}]=\CC$: у $\R$ ровно одно нетривиальное абелево расширение экспоненты $2$,
что согласуется с тем, что $\CC$~--- алгебраическое замыкание $\R$ (\rf{5.7}). (Поле $\R$ содержит $\mu_n$ лишь при $n\le2$.)
Для комплексного сопряжения $c$ имеем $\chi_{-1}(c)=c(i)/i=-1$.

(b) $F=\Q$, $n=2$. Каждое $a\in\Q^\times$ однозначно записывается в виде $a=\pm\prod_pp^{e_p}$ (почти все $e_p=0$), и $a$~---
квадрат в $\Q$ тогда и только тогда, когда знак~--- плюс и все $e_p$ чётны. Поэтому отображение
$a\mapsto\bigl(\text{знак},(e_p\bmod2)_p\bigr)$ задаёт изоморфизм
\[
\Q^\times/\Q^{\times2}\cong\F_2\oplus\bigoplus_{p\ \text{простое}}\F_2:
\]
это векторное пространство над $\F_2$ с базисом из классов $-1,2,3,5,7,\dots$; в каждом классе ровно одно целое число без
квадратных множителей. По теореме \rf{5.30} конечные абелевы расширения $\Q$ экспоненты, делящей $2$ (\emph{мультиквадратичные
поля} $\Q[\sqrt{d_1},\dots,\sqrt{d_r}]$), взаимно однозначно соответствуют конечномерным подпространствам этого
пространства. По \rf{5.30j}
\[
[\Q[\sqrt{d_1},\dots,\sqrt{d_r}]:\Q]=2^s,\qquad s=\dim\langle\bar d_1,\dots,\bar d_r\rangle,
\]
в частности степень равна $2^r$ тогда и только тогда, когда классы $d_1,\dots,d_r$ линейно независимы. Подпространству,
порождённому простыми $p_1,\dots,p_r$ (и $-1$), соответствует поле, полученное присоединением $\sqrt{p_1},\dots,\sqrt{p_r}$
(и $\sqrt{-1}$). По \rf{5.30k} квадратичные подполя поля $\Q[\sqrt{d_1},\dots,\sqrt{d_r}]$~--- это $\Q[\sqrt d]$, где $\bar d$
пробегает $2^s-1$ ненулевых векторов подпространства $\langle\bar d_1,\dots,\bar d_r\rangle$.
\end{example}

\begin{example}[5.31a: мультиквадратичные поля]
(1) $E=\Q[\sqrt2,\sqrt3]$. Классы $2$ и $3$ независимы, поэтому $[E:\Q]=4$ и $B(E)/\Q^{\times2}=\{1,2,3,6\}$ (по
представителям). Вложение \rf{5.30f}(c) $G\hookrightarrow\mu_2\times\mu_2$ биективно (порядки равны $4$), так что
$G=\{1,\sigma,\rho,\sigma\rho\}$, где $\sigma$ меняет знак у $\sqrt2$ и не двигает $\sqrt3$, а $\rho$~--- наоборот.
Изоморфизм $\kappa_E$ задаётся таблицей значений $\chi_a(g)=g(\sqrt a)/\sqrt a$:
\begin{center}
\begin{tabular}{c|cccc}
 & $1$ & $\sigma$ & $\rho$ & $\sigma\rho$\\\hline
$\chi_2$ & $1$ & $-1$ & $\phantom-1$ & $-1$\\
$\chi_3$ & $1$ & $\phantom-1$ & $-1$ & $-1$\\
$\chi_6$ & $1$ & $-1$ & $-1$ & $\phantom-1$
\end{tabular}
\end{center}
В каждой строке и в каждом столбце, кроме первого, есть $-1$: так выглядит невырожденность. Квадратичные подполя~---
$\Q[\sqrt2]$, $\Q[\sqrt3]$, $\Q[\sqrt6]$, и других нет; $\sqrt a\in E\iff a\in\{1,2,3,6\}\cdot\Q^{\times2}$. Например,
$\sqrt{24}=2\sqrt6\in E$, но $\sqrt5,\sqrt{-1}\notin E$. Поле $\Q[\sqrt2,\sqrt3,\sqrt6]$ совпадает с $E$ и имеет степень
$4$, а не $8$: $6=2\cdot3$. Элемент $\sqrt2+\sqrt3$ порождает $E$; его минимальный многочлен~--- $X^4-10X^2+1$.

\begin{keepfig}
\begin{center}
\begin{tikzpicture}[every node/.style={font=\small}]
\node (E) at (0,2.4) {$\Q[\sqrt2,\sqrt3]$};
\node (A) at (-1.9,1.2) {$\Q[\sqrt2]$};
\node (B) at (0,1.2) {$\Q[\sqrt3]$};
\node (C) at (1.9,1.2) {$\Q[\sqrt6]$};
\node (Q) at (0,0) {$\Q$};
\foreach \x in {A,B,C} {\draw[cold] (Q)--(\x); \draw[cold] (\x)--(E);}
\begin{scope}[xshift=7.2cm]
\node[cnew] (E1) at (0,2.4) {$\{1,2,3,6\}$};
\node[cnew] (A1) at (-1.9,1.2) {$\{1,2\}$};
\node[cnew] (B1) at (0,1.2) {$\{1,3\}$};
\node[cnew] (C1) at (1.9,1.2) {$\{1,6\}$};
\node[cnew] (Q1) at (0,0) {$\{1\}$};
\foreach \x in {A1,B1,C1} {\draw[cnew] (Q1)--(\x); \draw[cnew] (\x)--(E1);}
\end{scope}
\draw[<->,thick] (2.9,1.2) -- node[above,font=\footnotesize]{$K\leftrightarrow B(K)$} (4.4,1.2);
\end{tikzpicture}
\end{center}
\figcap{Слева~--- подполя $\Q[\sqrt2,\sqrt3]$, справа~--- соответствующие подгруппы $B(K)/\Q^{\times2}$ (по представителям).
Соответствие сохраняет включения и степени: $[K:\Q]=(B(K):\Q^{\times2})$.}
\end{keepfig}

(2) $E=\Q[\sqrt6,\sqrt{10},\sqrt{15}]$. В координатах по базису $\bar2,\bar3,\bar5$ классы $6,10,15$~--- это
$(1,1,0)$, $(1,0,1)$, $(0,1,1)$; их сумма равна нулю ($6\cdot10\cdot15=30^2$), а любые два независимы. Значит,
$[E:\Q]=4$, $E=\Q[\sqrt6,\sqrt{10}]\ni\sqrt{15}=\sqrt6\sqrt{10}/2$, и квадратичные подполя~--- $\Q[\sqrt6]$,
$\Q[\sqrt{10}]$, $\Q[\sqrt{15}]$.

(3) Для различных простых $p_1,\dots,p_k$ поле $\Q[\sqrt{p_1},\dots,\sqrt{p_k}]$ имеет степень $2^k$ и группу
$(\Z/2)^k$; его квадратичные подполя~--- $\Q[\sqrt{p_S}]$, где $p_S=\prod_{i\in S}p_i$, $\varnothing\ne
S\subset\{1,\dots,k\}$. В башне $\Q\subset\Q[\sqrt{p_1}]\subset\dots\subset\Q[\sqrt{p_1},\dots,\sqrt{p_k}]$ каждый шаг
имеет степень $2$ с базисом $\{1,\sqrt{p_j}\}$, поэтому $2^k$ чисел $\sqrt{p_S}$ (включая $\sqrt{p_\varnothing}=1$)
образуют базис поля над $\Q$ (\rf{1.20}); в частности, они линейно независимы над $\Q$. Например, при $k=3$ поле
$\Q[\sqrt2,\sqrt3,\sqrt5]$ имеет степень $8$ и семь квадратичных подполей $\Q[\sqrt d]$, $d\in\{2,3,5,6,10,15,30\}$.
\end{example}

\begin{example}[5.31b: квадратные корни в круговых полях]
Применим \rf{5.30m} с $n=2$ к $E=\Q[\zeta_m]$, $\zeta_m=e^{2\pi i/m}$, $G\cong(\Z/m)^\times$ (\rf{5.11a}). Наибольшее подполе $E$
экспоненты, делящей $2$, имеет степень $|\Hom(G,\mu_2)|=(G:G^2)$; все квадратичные подполя $E$ лежат в нём, и по \rf{5.30k} их
ровно $(G:G^2)-1$. При $m=8$ и простом $m$ сами эти подполя найдены в разделе о круговых расширениях (\rf{5.11g}, \rf{5.11f});
теория Куммера объясняет их число.
\begin{itemize}[leftmargin=*]
\item $m=8$: все элементы $(\Z/8)^\times=\{1,3,5,7\}$ имеют квадрат $1$, так что само $E$ имеет экспоненту $2$ и
$(B(E):\Q^{\times2})=4$. Так как $i=\zeta_8^2\in E$ и $\sqrt2=\zeta_8+\zeta_8^{-1}\in E$, а классы $-1$, $2$ независимы,
получаем $B(E)=\langle-1,2\rangle\Q^{\times2}$, т.\,е.\ $\Q[\zeta_8]=\Q[i,\sqrt2]$ с тремя квадратичными подполями
$\Q[i]$, $\Q[\sqrt2]$, $\Q[\sqrt{-2}]$ (ср.\ \rf{5.11g}).
\item $m=12$ и $m=24$: группы $(\Z/12)^\times$ и $(\Z/24)^\times$ (порядков $4$ и $8$) тоже имеют экспоненту $2$. Так как
$i=\zeta_{24}^6$, $\sqrt2=\zeta_{24}^3+\zeta_{24}^{-3}$, $\sqrt3=\zeta_{24}^2+\zeta_{24}^{-2}$, получаем
$\Q[\zeta_{12}]=\Q[i,\sqrt3]$ и $\Q[\zeta_{24}]=\Q[i,\sqrt2,\sqrt3]$ (сравните степени: $2^3=8=\varphi(24)$).
\item $m=p$ нечётное простое: $G$ циклична (\rf{5.7c}) чётного порядка $p-1$, так что $(G:G^2)=2$. Значит, в $\Q[\zeta_p]$
ровно одно квадратичное подполе, и $\sqrt a\in\Q[\zeta_p]$ ровно для одного нетривиального класса
$a\in\Q^\times/\Q^{\times2}$, а именно для класса $p^*=(-1)^{(p-1)/2}p$: сумма Гаусса
$g=\sum_{k=1}^{p-1}\bigl(\frac kp\bigr)\zeta_p^k$ удовлетворяет $g^2=p^*$ (\rf{5.11f}). Например, $\sqrt{-3}=2\omega+1$
($\omega=\zeta_3$) при $p=3$; $\sqrt5=\zeta_5+\zeta_5^4-\zeta_5^2-\zeta_5^3$ при $p=5$;
$\sqrt{-7}=\zeta_7+\zeta_7^2+\zeta_7^4-\zeta_7^3-\zeta_7^5-\zeta_7^6$ при $p=7$ (\rf{5.11h}, ср.\ \rf{3.22});
$\sqrt{13}=\sum_{k=1}^{12}\bigl(\frac k{13}\bigr)\zeta_{13}^k$ при $p=13$. Поэтому, например,
$\sqrt2,\sqrt{-13}\notin\Q[\zeta_{13}]$.
\end{itemize}
\end{example}

\begin{example}[5.31c: кубические корни над {$\Q[\omega]$}]
Пусть $F=\Q[\omega]$, $\omega=e^{2\pi i/3}$, $n=3$; $\mu_3=\{1,\omega,\omega^2\}\subset F$.

\emph{Независимость классов.} Комплексное сопряжение~--- нетривиальный автоморфизм $F$, поэтому $N(x)=x\bar x$~---
гомоморфизм $F^\times\to\Q^\times$ (значения лежат в $F\cap\R=\Q$; это норма, ср.\ \rf{5.24}), и $N(q)=q^2$ при $q\in\Q^\times$. Если $2^a3^b=c^3$ с $c\in
F^\times$, то $2^{2a}3^{2b}=N(c)^3$~--- куб рационального числа, откуда $3\mid2a$ и $3\mid2b$, т.\,е.\ $3\mid a$ и $3\mid
b$. Значит, классы $2$ и $3$ порождают в $F^\times/F^{\times3}$ подгруппу, изоморфную $(\Z/3)^2$. (Так же классы любых
различных простых чисел независимы.)

(1) Класс $2$ имеет порядок $3$, поэтому $X^3-2$ неприводим над $F$, и $F[\sqrt[3]2]/F$ циклично степени $3$ с
образующей $\sigma\colon\sqrt[3]2\mapsto\omega\sqrt[3]2$. Далее, $F[\sqrt[3]2]=F[\sqrt[3]4]$, так как $4=2^2$ порождает ту
же подгруппу (\rf{5.30j}(c); ср.\ пример после \rf{5.28}); явно, $\sqrt[3]4=(\sqrt[3]2)^2$ и $\sqrt[3]2=(\sqrt[3]4)^2/2$.

(2) $E=F[\sqrt[3]2,\sqrt[3]3]$: $B(E)/F^{\times3}=\langle\bar2,\bar3\rangle\cong(\Z/3)^2$, поэтому $[E:F]=9$, и
$G\cong(\Z/3)^2$ порождается автоморфизмами
\[
\sigma\colon\sqrt[3]2\mapsto\omega\sqrt[3]2,\ \sqrt[3]3\mapsto\sqrt[3]3;\qquad
\tau\colon\sqrt[3]2\mapsto\sqrt[3]2,\ \sqrt[3]3\mapsto\omega\sqrt[3]3
\]
(вложение \rf{5.30f}(c) $G\hookrightarrow\mu_3^2$ биективно). Изоморфизм $\kappa_E$:
$\chi_{2^x3^y}(\sigma^i\tau^j)=\omega^{xi+yj}$~--- <<скалярное произведение>> на $(\Z/3)^2$. В $(\Z/3)^2$ ровно четыре
подгруппы порядка $3$: $\langle\bar2\rangle$, $\langle\bar3\rangle$, $\langle\bar6\rangle$, $\langle\overline{12}\rangle$
(класс $18=2\cdot3^2$ лежит в $\langle\overline{12}\rangle$: $18^2=12\cdot27$). По \rf{5.30k} у $E/F$ ровно четыре промежуточных
поля степени $3$:
\[
F[\sqrt[3]2],\quad F[\sqrt[3]3],\quad F[\sqrt[3]6],\quad F[\sqrt[3]{12}].
\]
Например, $\Gal(E/F[\sqrt[3]6])=\{\sigma^i\tau^j\mid i+j\equiv0\pmod 3\}=\langle\sigma\tau^{-1}\rangle$: действительно,
$\sigma\tau^{-1}(\sqrt[3]6)=\omega\cdot\omega^{-1}\sqrt[3]6=\sqrt[3]6$.

\begin{keepfig}
\begin{center}
\begin{tikzpicture}[every node/.style={font=\small}]
\node (E) at (0,2.5) {$E=F[\sqrt[3]2,\sqrt[3]3]$};
\node (A) at (-2.7,1.25) {$F[\sqrt[3]2]$};
\node (B) at (-0.9,1.25) {$F[\sqrt[3]3]$};
\node (C) at (0.9,1.25) {$F[\sqrt[3]6]$};
\node (D) at (2.7,1.25) {$F[\sqrt[3]{12}]$};
\node (F) at (0,0) {$F=\Q[\omega]$};
\foreach \x in {B,C,D} {\draw[cold] (F)--(\x); \draw[cold] (\x)--(E);}
\draw[cold] (F) -- node[below left,font=\footnotesize]{$3$} (A);
\draw[cold] (A) -- node[above left,font=\footnotesize]{$3$} (E);
\begin{scope}[xshift=8.6cm]
\node[cnew] (E1) at (0,2.5) {$\langle\bar2,\bar3\rangle\cong(\Z/3)^2$};
\node[cnew] (A1) at (-2.4,1.25) {$\langle\bar2\rangle$};
\node[cnew] (B1) at (-0.8,1.25) {$\langle\bar3\rangle$};
\node[cnew] (C1) at (0.8,1.25) {$\langle\bar6\rangle$};
\node[cnew] (D1) at (2.4,1.25) {$\langle\overline{12}\rangle$};
\node[cnew] (F1) at (0,0) {$1$};
\foreach \x in {A1,B1,C1,D1} {\draw[cnew] (F1)--(\x); \draw[cnew] (\x)--(E1);}
\end{scope}
\draw[<->,thick] (3.9,1.25) -- (5.6,1.25);
\end{tikzpicture}
\end{center}
\figcap{Подполя $E=\Q[\omega,\sqrt[3]2,\sqrt[3]3]$ над $F=\Q[\omega]$ (каждый шаг степени $3$) и подгруппы
$B(E)/F^{\times3}\cong(\Z/3)^2$~--- четыре прямые в плоскости над $\F_3$.}
\end{keepfig}

(3) \emph{Следствия над $\Q$.} $\sqrt[3]3\notin\Q[\sqrt[3]2]$: иначе $\sqrt[3]3\in F[\sqrt[3]2]$, т.\,е.\
$3\in\langle2\rangle F^{\times3}$ по \rf{5.30j}(b). Поле разложения $\Q[\omega,\sqrt[3]2,\sqrt[3]3]$ многочлена
$(X^3-2)(X^3-3)$ имеет степень $2\cdot9=18$ над $\Q$, а вещественное поле $\Q[\sqrt[3]2,\sqrt[3]3]$ (не содержащее
$\omega$)~--- степень $9$. Более того, $\theta=\sqrt[3]2+\sqrt[3]3$ имеет над $\Q$ степень $9$, т.\,е.\ порождает
$\Q[\sqrt[3]2,\sqrt[3]3]$. Действительно, образы $\sigma^i\tau^j(\theta)=\omega^i\sqrt[3]2+\omega^j\sqrt[3]3$ попарно
различны: если $\omega^i\sqrt[3]2+\omega^j\sqrt[3]3=\omega^k\sqrt[3]2+\omega^l\sqrt[3]3$ при $(i,j)\ne(k,l)$, то $i\ne k$ и
$j\ne l$ (иначе равенство сразу даёт $(i,j)=(k,l)$), и
$\sqrt[3]3/\sqrt[3]2=(\omega^i-\omega^k)/(\omega^l-\omega^j)\in F$, т.\,е.\ $3\cdot2^{-1}\in F^{\times3}$, что
противоречит независимости классов $2$ и $3$. Поэтому стабилизатор $\theta$ в $G$ тривиален и $F[\theta]=E$,
$[F[\theta]:F]=9$; так как $F[\theta]=\Q[\theta][\omega]$ и $[\Q[\theta][\omega]:\Q[\theta]]\le2$, получаем
$[\Q[\theta]:\Q]\ge18/2=9$. Но $\theta\in\Q[\sqrt[3]2,\sqrt[3]3]$, поэтому $\Q[\theta]=\Q[\sqrt[3]2,\sqrt[3]3]$.
(Минимальный многочлен $\theta$ над $\Q$~--- $X^9-15X^6-87X^3-125$.)
\end{example}

\begin{example}[5.31d: корни четвёртой степени над {$\Q[i]$}]
Пусть $F=\Q[i]$, $n=4$, $\mu_4=\{\pm1,\pm i\}\subset F$. Порядок класса в $F^\times/F^{\times4}$ равен $1$, $2$ или $4$, и
по \rf{5.30i}(b) $X^4-a$ неприводим над $F$ тогда и только тогда, когда $a$ не квадрат в $F$. По \rf{5.30j}(b) (над $\Q$, $n=2$)
рациональное $d$ является квадратом в $\Q[i]$ лишь при $d\in\Q^{\times2}\cup(-\Q^{\times2})$; поэтому $\pm2$, $\pm3$~---
не квадраты в $F$, а $4$, $9$, $-1=i^2$, $-4=(2i)^2$~--- квадраты.
\begin{itemize}[leftmargin=*]
\item $a=2$: $X^4-2$ неприводим над $F$, и $E=F[\sqrt[4]2]/F$ циклично степени $4$ с образующей
$\sigma\colon\sqrt[4]2\mapsto i\sqrt[4]2$; $\chi_{2^k}(\sigma^j)=i^{kj}$. Единственное промежуточное поле~--- $F[\sqrt[4]4]=
F[\sqrt2]=\Q[\zeta_8]$; его неподвижная группа $\langle\sigma^2\rangle$: $\sigma^2(\sqrt2)=(\sigma^2\sqrt[4]2)^2=
(-\sqrt[4]2)^2=\sqrt2$. Над $\Q$ поле $E=\Q[i,\sqrt[4]2]$~--- поле разложения $X^4-2$, степени $8$.
\item $a=4$: класс порядка $2$ (если $c^4=4$, то $c^2=\pm2$~--- невозможно). Поэтому $4$ не является четвёртой
степенью, и всё же $X^4-4=(X^2-2)(X^2+2)$, $[F[\sqrt[4]4]:F]=2$.
\item $a=-4=(1+i)^4$: класс тривиален, и $X^4+4=\prod_{k=0}^{3}\bigl(X-i^k(1+i)\bigr)$ раскладывается над $F$ на линейные
множители. (Над $\Q$: $X^4+4=(X^2+2X+2)(X^2-2X+2)$, ср.\ \rf{5.29}(a).)
\item $a=-1$: класс порядка $2$, $X^4+1=(X^2-i)(X^2+i)$, и $F[\sqrt[4]{-1}]=F[\zeta_8]=\Q[\zeta_8]$.
\item \emph{Экспонента делит $n$, но не равна $n$.} Класс $9$ имеет порядок $2$ ($9=3^2$, а $\pm3$~--- не квадраты в $F$).
Корни степени $4$ из $9$~--- это $\pm\sqrt3$, $\pm i\sqrt3$, поэтому подгруппе $B=\langle9\rangle F^{\times4}$
соответствует $F[B^{1/4}]=\Q[i,\sqrt3]$ с группой $\Z/2$ экспоненты $2\ne4$; при этом $X^4-9=(X^2-3)(X^2+3)$.
\end{itemize}
\end{example}

\begin{example}[5.31e: {$\Q[\zeta_5]$}, $n=5$]
Пусть $F=\Q[\zeta_5]$; $[F:\Q]=4$. Класс $2$ в $F^\times/F^{\times5}$ нетривиален: иначе какой-то корень пятой степени
из $2$, а с ним (после умножения на элемент $\mu_5\subset F$) и вещественный $\sqrt[5]2$, лежал бы в $F$, но
$[\Q[\sqrt[5]2]:\Q]=5$ не делит $4$. Порядок класса~--- простое число $5$, поэтому $X^5-2$ неприводим над $F$ (\rf{5.30i}(c)),
$F[\sqrt[5]2]/F$ циклично степени $5$, а $[F[\sqrt[5]2]:\Q]=20$ (ср.\ \rf{3.23}). Поле $F[\sqrt[5]2,\sqrt[5]4,\sqrt[5]8]$
совпадает с $F[\sqrt[5]2]$: все три класса лежат в $\langle\bar2\rangle$.
\end{example}

\begin{example}[5.31f: конечные поля]
Пусть $F=\F_q$ и $n\mid q-1$ (это и есть условие $\mu_n\subset\F_q$: группа $\F_q^\times$ циклична порядка $q-1$, \rf{5.7c}).
Возведение в степень $n$ на циклической группе $\F_q^\times$ имеет ядро $\mu_n$ порядка $n$, поэтому
$\F_q^\times/\F_q^{\times n}\cong\Z/n$; его порождает класс образующей $g$ группы $\F_q^\times$. В $\Z/n$ ровно по одной
подгруппе порядка $d$ для каждого $d\mid n$, она порождена классом $g^{n/d}$. По теореме \rf{5.30} для каждого $d\mid n$ у
$\F_q$ ровно одно расширение степени $d$, абелево с экспонентой, делящей $n$, а именно
$\F_q[(g^{n/d})^{1/n}]=\F_q[g^{1/d}]$; конечно, это $\F_{q^d}$ (\rf{4.23}), и теория Куммера воспроизводит известное.

\emph{Изоморфизм $\kappa_E$ и критерий Эйлера.} Пусть $E=\F_{q^n}=\F_q[g^{1/n}]$. По теореме \rf{5.30}
$B(E)=\langle g\rangle\F_q^{\times n}=\F_q^\times$: каждый элемент $\F_q$ является $n$-й степенью в $\F_{q^n}$. Группа
$\Gal(E/\F_q)$ порождена автоморфизмом Фробениуса $\Frob(x)=x^q$ (\rf{4.23}). Если $a\in\F_q^\times$ и $\alpha\in E$, $\alpha^n=a$, то
\[
\chi_a(\Frob)=\frac{\alpha^q}{\alpha}=\alpha^{q-1}=a^{(q-1)/n}.
\]
Так как характер $\chi_a$ определяется значением на образующей, а ядро $a\mapsto\chi_a$ равно $\F_q^{\times n}$ (\rf{5.30d}),
получаем \emph{критерий Эйлера}: $a\in\F_q^{\times n}\iff a^{(q-1)/n}=1$. (При $q=p$, $n=2$ это критерий Эйлера,
использованный в доказательстве \rf{5.11l}: $a^{(p-1)/2}\bmod p$~--- символ Лежандра $\bigl(\frac ap\bigr)$ из \rf{5.11f}.)

\emph{Численно.} $q=7$, $n=3$: кубы в $\F_7^\times$~--- это $\{1,6\}$, и $\F_7^\times/\F_7^{\times3}$ состоит из классов
$\{1,6\}$, $\{2,5\}$, $\{3,4\}$. По \rf{5.30i}(c) $X^3-a$ неприводим над $\F_7$ при $a\in\{2,3,4,5\}$ и раскладывается на
линейные множители при $a\in\{1,6\}$. Так как $2\cdot3=6$~--- куб, класс $3$ равен $\bar2^{-1}$, и
$\F_{343}=\F_7[\sqrt[3]2]=\F_7[\sqrt[3]3]$. Критерий Эйлера: $2^{(7-1)/3}=4\ne1$, а $6^2=1$. При $n=6$ имеем
$\F_7^{\times6}=\{1\}$, $\F_7^\times\cong\Z/6$ с образующей $3$, и подполя $\F_{7^6}=\F_7[\sqrt[6]3]$~--- это $\F_7$,
$\F_{49}=\F_7[\sqrt3]$, $\F_{343}=\F_7[\sqrt[3]3]$ и $\F_{7^6}$; класс $2$ имеет порядок $3$, поэтому
$X^6-2=(X^3-3)(X^3+3)$ над $\F_7$.
\end{example}

\begin{example}[5.31g: поле $\CC(t)$]
Пусть $F=\CC(t)$, $n\ge1$ любое. Кольцо $\CC[t]$ факториально, его простые элементы (с точностью до множителей из
$\CC^\times$)~--- это $t-c$, $c\in\CC$, поэтому каждое $f\in\CC(t)^\times$ однозначно записывается как
$f=u\prod_c(t-c)^{m_c}$, $u\in\CC^\times$, $m_c\in\Z$ (почти все $m_c=0$). Так как $\CC^\times=\CC^{\times n}$, имеем
$f\in\CC(t)^{\times n}\iff n\mid m_c$ для всех $c$. Значит,
\[
\CC(t)^\times/\CC(t)^{\times n}\cong\bigoplus_{c\in\CC}\Z/n,\qquad f\mapsto(m_c\bmod n)_c,
\]
~--- свободный $\Z/n$-модуль с базисом из классов $t-c$. Следствия:
\begin{itemize}[leftmargin=*]
\item для различных $c_1,\dots,c_k$ поле $\CC(t)[(t-c_1)^{1/n},\dots,(t-c_k)^{1/n}]$ имеет степень $n^k$ и группу
$(\Z/n)^k$; в частности, $X^n-t$ неприводим, и $\CC(t)[t^{1/n}]/\CC(t)$ циклично степени $n$;
\item при $n=2$ квадратичные расширения $\CC(t)$~--- это $\CC(t)\bigl[\sqrt{(t-c_1)\cdots(t-c_k)}\bigr]$ для конечных
непустых множеств $\{c_1,\dots,c_k\}\subset\CC$ различных чисел, и два таких поля совпадают тогда и только тогда, когда
совпадают множества;
\item например, $\CC(t)[\sqrt t,\sqrt{t-1}]$ имеет степень $4$ и ровно три квадратичных подполя: $\CC(t)[\sqrt t]$,
$\CC(t)[\sqrt{t-1}]$, $\CC(t)[\sqrt{t(t-1)}]$.
\end{itemize}
\end{example}

\begin{example}[5.31h: поле $\Qp$, $n=2$]
Этот пример опирается на гл.~7: $p$-адические числа и лемма Гензеля вводятся в \rf{7.33}--\rf{7.33a}. Пусть $p$ нечётно.
Каждый $x\in\Qp^\times$ однозначно записывается как $x=p^vu$, $v\in\Z$, $u\in\Zp^\times$ (\rf{7.33}). По лемме Гензеля \rf{7.33a}, применённой к $X^2-u$ (производная $2a_0\not\equiv0\pmod p$), единица $u$ является квадратом в $\Zp$
тогда и только тогда, когда её вычет по модулю $p$~--- квадрат в $\F_p$. Если $x=y^2$, $y=p^wu'$, то $v=2w$ и $u=u'^2$.
Поэтому $x\in\Qp^{\times2}$ $\iff$ $v$ чётно и вычет $u$~--- квадрат, и
\[
\Qp^\times/\Qp^{\times2}=\{1,\ u_0,\ p,\ u_0p\}\cong(\Z/2)^2,
\]
где $u_0\in\Z$~--- невычет по модулю $p$. По теореме \rf{5.30} у $\Qp$ ровно три квадратичных расширения $\Qp[\sqrt{u_0}]$,
$\Qp[\sqrt p]$, $\Qp[\sqrt{u_0p}]$, а их композит $\Qp[\sqrt{u_0},\sqrt p]$ степени $4$~--- наибольшее абелево расширение
экспоненты $2$. Например, у $\Q_3$ это $\Q_3[\sqrt{-1}]$, $\Q_3[\sqrt3]$, $\Q_3[\sqrt{-3}]$, а у $\Q_5$~--- $\Q_5[\sqrt2]$,
$\Q_5[\sqrt5]$, $\Q_5[\sqrt{10}]$ (при этом $\sqrt{-1}\in\Q_5$, так как $-1\equiv2^2\pmod5$).

При $p=2$ по \rf{7.33b}(b) $\Z_2^\times=\{\pm1\}\times(1+4\Z_2)$ и $1+4\Z_2=\{5^a\mid a\in\Z_2\}\cong\Z_2$; квадраты единиц~---
это $5^{2\Z_2}$, подгруппа индекса $2$ в $1+4\Z_2$. Поэтому $\Z_2^\times/\Z_2^{\times2}$ имеет представители $\pm1,\pm5$,
группа $\Q_2^\times/\Q_2^{\times2}$ имеет порядок $8$ (представители $\pm1,\pm5,\pm2,\pm10$), и у $\Q_2$ ровно семь
квадратичных расширений.
\end{example}

\begin{example}[5.31i: без корней из единицы]
В этом примере условие $\zeta\in F$ не предполагается.

(1) Поле $\Q[\sqrt[3]2]$ не является расширением Галуа над $\Q$: оно вещественно, а другие корни $\omega\sqrt[3]2$,
$\omega^2\sqrt[3]2$ неприводимого многочлена $X^3-2$ не вещественны. Причина~--- $\mu_3\not\subset\Q$. Вообще, если
$\alpha^n=a\in F^\times$ и $\operatorname{char}F\nmid n$, то $F[\alpha]$~--- поле разложения $X^n-a$ тогда и только тогда,
когда $F[\alpha]$ содержит первообразный корень $\eta\in\Om$ степени $n$ из $1$ (ср.\ \rf{2.10}): если $F[\alpha]$ содержит
все корни $X^n-a$, то $\eta\alpha\in F[\alpha]$ и $\eta=\eta\alpha/\alpha$, а обратное очевидно, так как корни~--- это
$\eta^i\alpha$. Для вещественного поля $F[\alpha]$ это возможно лишь при $n\le2$.

\emph{Как помогает присоединение $\mu_n$.} Как отмечено в \rf{5.27a}, условие обходят, присоединяя сначала корни из единицы.
Для $X^3-2$: $\Q[\omega]/\Q$~--- круговой шаг степени $2$ (\rf{5.8}), над $\Q[\omega]$ класс $2$ имеет порядок $3$ (\rf{5.31c}), и
$\Q[\omega,\sqrt[3]2]/\Q[\omega]$~--- циклическое расширение Куммера степени $3$ с группой, изоморфной $\mu_3$
($\sigma\mapsto\sigma(\sqrt[3]2)/\sqrt[3]2$). Поле $\Q[\omega,\sqrt[3]2]$~--- поле разложения $X^3-2$ над $\Q$, степени $6$,
с группой $S_3$; шаг Куммера соответствует нормальной подгруппе $A_3$.
\begin{keepfig}
\begin{center}
\begin{tikzpicture}[every node/.style={font=\small}]
\node (T) at (0,2.7) {$\Q[\omega,\sqrt[3]2]$};
\node (A) at (-3.4,1.25) {$\Q[\sqrt[3]2]$};
\node (B) at (-1.6,1.25) {$\Q[\omega\sqrt[3]2]$};
\node (C) at (0.3,1.25) {$\Q[\omega^2\sqrt[3]2]$};
\node (W) at (2.9,1.25) {$\Q[\omega]$};
\node (Q) at (0,0) {$\Q$};
\foreach \x in {A,B,C} {\draw (Q)--(\x); \draw (\x)--(T);}
\draw[cold,thick] (Q) -- node[below right,font=\footnotesize]{$2$, круговой шаг} (W);
\draw[cnew,thick] (W) -- node[above right,font=\footnotesize]{$3$, шаг Куммера ($\mu_3\subset\Q[\omega]$)} (T);
\end{tikzpicture}
\end{center}
\figcap{Подполя поля разложения $X^3-2$ над $\Q$. Три кубических подполя (слева) не нормальны над $\Q$; после
присоединения $\mu_3$ остаётся циклическое расширение Куммера (красное).}
\end{keepfig}

(2) \emph{Циклическое расширение без радикала.} Поле $K=\Q[\cos\frac{2\pi}7]\subset\Q[\zeta_7]$ циклично степени $3$ над
$\Q$ (\rf{5.11d}, \rf{5.11h}; ср.\ \rf{3.22}), но, как показано в \rf{5.27a}, не порождено над $\Q$ кубическим корнем; так что без условия
$\mu_n\subset F$ утверждение \rf{5.30l} неверно. После присоединения $\omega$ всё в порядке: $K\cap\Q[\omega]=\Q$ (степени $3$
и $2$), так что по \rf{3.19} $K[\omega]/\Q[\omega]$ циклично степени $3$ и по \rf{5.30l} порождено кубическим корнем. Найдём его
явно. Пусть $\theta_0=2\cos\frac{2\pi}7=\zeta_7+\zeta_7^{-1}$; автоморфизм $\sigma\colon\zeta_7\mapsto\zeta_7^2$ порождает
$\Gal(K/\Q)$ (его порядок на $K$ равен $3$), продолжается до образующей $\Gal(K[\omega]/\Q[\omega])$ с
$\sigma(\omega)=\omega$ (\rf{3.19}) и даёт сопряжённые $\theta_1=\sigma\theta_0=\theta_0^2-2$, $\theta_2=\sigma\theta_1=
\theta_1^2-2$~--- корни $X^3+X^2-2X-1$ (\rf{5.11h}). Резольвента Лагранжа $r=\theta_0+\omega\theta_1+\omega^2\theta_2$
удовлетворяет $\sigma(r)=\omega^2r$, поэтому $r^3\in\Q[\omega]$.

\emph{Вычисление $r^3$.} Так как $\sigma(\theta_i)=\theta_i^2-2=\theta_{i+1}$ (индексы по модулю $3$), имеем
$\theta_i^2=\theta_{i+1}+2$, а коэффициенты многочлена $X^3+X^2-2X-1$ дают
\[
\theta_0+\theta_1+\theta_2=-1,\qquad\theta_0\theta_1+\theta_1\theta_2+\theta_2\theta_0=-2,\qquad\theta_0\theta_1\theta_2=1.
\]
Отсюда (все суммы~--- по $i\in\Z/3$)
\begin{gather*}
\textstyle\sum_i\theta_i^3=\sum_i\theta_i\theta_{i+1}+2\sum_i\theta_i=-4,\qquad
\sum_i\theta_i\theta_{i+1}^2=\sum_i\theta_i\theta_{i+2}+2\sum_i\theta_i=-4,\\
\textstyle\sum_i\theta_i^2\theta_{i+1}=\sum_i\bigl(\theta_{i+1}^2+2\theta_{i+1}\bigr)=
\sum_i\bigl(\theta_{i+2}+2+2\theta_{i+1}\bigr)=3.
\end{gather*}
Раскрыв куб и учтя $\omega^3=1$, для $r=a+\omega b+\omega^2c$ получаем
\[
r^3=a^3+b^3+c^3+6abc+3\omega(a^2b+b^2c+c^2a)+3\omega^2(ab^2+bc^2+ca^2).
\]
При $a=\theta_0$, $b=\theta_1$, $c=\theta_2$ с учётом $\omega^2=-1-\omega$ это даёт
\[
r^3=-4+6+9\omega-12\omega^2=14+21\omega=7(2+3\omega),\qquad\text{т.\,е.}\qquad
K[\omega]=\Q[\omega]\bigl[\sqrt[3]{7(2+3\omega)}\bigr]
\]
(так как $r^3\ne0$, имеем $\sigma(r)=\omega^2r\ne r$, т.\,е.\ $r\notin\Q[\omega]$, а степень $[K[\omega]:\Q[\omega]]=3$
проста). Далее, $r'=\theta_0+\omega^2\theta_1+\omega\theta_2=\bar r$ и $rr'=|r|^2=7$ (так как
$|r|^6=|14+21\omega|^2=14^2-14\cdot21+21^2=7^3$), а $(\theta_0+\theta_1+\theta_2)+r+r'=3\theta_0$, так как
$1+\omega+\omega^2=0$. Отсюда формула типа Кардано:
\[
2\cos\frac{2\pi}7=\frac13\Bigl(-1+r+\frac7r\Bigr),\qquad r^3=14+21\omega;
\]
если подставить в неё вместо $r$ два других кубических корня $\omega r$, $\omega^2r$ из $14+21\omega$, получатся
сопряжённые $\theta_2$ и $\theta_1$.

(3) При $n=2$ и $\operatorname{char}F\ne2$ условие $\mu_2=\{\pm1\}\subset F$ выполнено всегда, так что теория Куммера
при $n=2$ применима к любому такому полю (\rf{5.27a}; ср.\ упражнение 2-1).
\end{example}

\begin{remark}[5.31j: характеристика, делящая $n$]
Пусть $\operatorname{char}F=p\mid n$. Тогда первообразных корней степени $n$ из $1$ нет, а радикалы степени $p$
бесполезны: многочлен $X^p-t$ над $\F_p(t)$ неприводим (\rf{2.18}), но $X^p-t=(X-t^{1/p})^p$, и поле $\F_p(t)[t^{1/p}]$,
порождённое его корнем, чисто несепарабельно над $\F_p(t)$. Вместо уравнения $X^p=a$ здесь работает уравнение
Артина--Шрайера $X^p-X=a$ (\rf{5.32c}), для степеней $p^k$~--- векторы Витта (\rf{5.32s}, \rf{5.32t}), а абелево расширение
экспоненты, делящей $n=p^km$ ($p\nmid m$), раскладывается в композит $p$-части и части экспоненты, делящей $m$, к
которой при $\mu_m\subset F$ применима теория Куммера (подробно~--- \rf{5.32x}).
\end{remark}

\begin{remark}[5.32: совершенное спаривание]
Пусть $E/F$~--- конечное абелево расширение, экспонента которого делит $n$, $G=\Gal(E/F)$ и
\[
B(E)=\{a\in F^\times\mid a\ \text{становится $n$-й степенью в}\ E\}.
\]
Тогда спаривание
\[
\langle a,\sigma\rangle=\frac{\sigma(a^{1/n})}{a^{1/n}}\colon\qquad B(E)/F^{\times n}\times G\to\mu_n
\]
совершенно. В частности, имеются канонические изоморфизмы
\[
B(E)/F^{\times n}\xrightarrow{\ \sim\ }\Hom(G,\mu_n),\qquad G\xrightarrow{\ \sim\ }\Hom(B(E)/F^{\times n},\mu_n),
\]
а неканонически $B(E)/F^{\times n}\cong G$ (\rf{5.30a}(b)). На языке подгрупп: для $B$ из множества (b) теоремы \rf{5.30} и
$E_B=F[B^{1/n}]$
\[
\Gal(E_B/F)\cong\Hom(B/F^{\times n},\mu_n),\ \ \sigma\mapsto\Bigl(b\mapsto\frac{\sigma(b^{1/n})}{b^{1/n}}\Bigr),\qquad
B/F^{\times n}\cong\Hom(\Gal(E_B/F),\mu_n).
\]
\end{remark}

\begin{proof}
\emph{Корректность.} Значение $\langle a,\sigma\rangle$ лежит в $\mu_n$ и не зависит от выбора корня $a^{1/n}\in E$
(\rf{5.30d}(a)). Оно зависит лишь от класса $a$ по модулю $F^{\times n}$: если $a'=ac^n$, $c\in F^\times$, то $a^{1/n}c$~---
корень из $a'$, и $\sigma(a^{1/n}c)/(a^{1/n}c)=\sigma(a^{1/n})/a^{1/n}$.

\emph{Билинейность.} Мультипликативность по $\sigma$~--- это гомоморфность $\chi_a$ (\rf{5.30d}(a)), по $a$~--- равенство
$\chi_{ab}=\chi_a\chi_b$ (\rf{5.30d}(b)).

\emph{Совершенность}~--- это \rf{5.30g} при $B=B(E)$ (по теореме \rf{5.30} $E=F[B(E)^{1/n}]$): левое ядро тривиально по
\rf{5.30d}(b), правое~--- так как автоморфизм, не двигающий все корни $a^{1/n}$, $a\in B(E)$, тождествен на $E$; обе группы
конечны, их экспоненты делят $n$, и остаётся применить \rf{5.30c}(a). Для $E=E_B$ имеем $B(E)=B$ по теореме \rf{5.30}.
\end{proof}

\paragraph{Соответствие Галуа как соответствие аннуляторов.} Для подгруппы $F^{\times n}\subset B'\subset B(E)$ её
аннулятор в $G$~--- это автоморфизмы, не двигающие корни $b^{1/n}$, $b\in B'$, т.\,е.\ $(B'/F^{\times n})^\perp=
\Gal(E/F[B'^{1/n}])$. Обратно, для $H\subset G$ имеем $H^\perp=B(E^H)/F^{\times n}$: $\langle a,\sigma\rangle=1$ для всех
$\sigma\in H$ $\iff$ $a^{1/n}\in E^H$. По \rf{5.30c}(b) отображение $B'\mapsto(B'/F^{\times n})^\perp$~--- обращающая
включения биекция, и $|\Gal(E/F[B'^{1/n}])|=|G|/(B':F^{\times n})$. Так для абелевых расширений экспоненты, делящей $n$,
основная теорема теории Галуа превращается в двойственность конечных абелевых групп (\rf{5.30c}), а поля нужны лишь для
построения спаривания (теорема Гильберта~90, т.\,е.\ \rf{5.23}, и равенство $E^G=F$).

\paragraph{Спаривание в примерах.} При $n=2$ группа $B(E)$ квадратичного расширения~--- это группа $S(E)$ из упражнения 2-1,
и совершенность означает, что $B(E)/F^{\times2}$ имеет порядок $2$: $B(E)=F^{\times2}\cup aF^{\times2}$, где
$E=F[\sqrt a]$. Таблица в \rf{5.31a}(1)~--- спаривание для $\Q[\sqrt2,\sqrt3]$; в \rf{5.31c}(2) спаривание~--- скалярное
произведение $\langle2^x3^y,\sigma^i\tau^j\rangle=\omega^{xi+yj}$ на $(\Z/3)^2$; в \rf{5.31d}~--- $\langle2^k,\sigma^j\rangle=
i^{kj}$ на $\Z/4$; для конечных полей (\rf{5.31f}) $\langle a,\Frob\rangle=a^{(q-1)/n}$, и невырожденность слева~--- критерий
Эйлера; для $\R$ (пример \rf{5.31}(a)) $\langle-1,c\rangle=-1$.

\paragraph{Бесконечный вариант.} Все абелевы расширения $F$ экспоненты, делящей $n$, можно рассмотреть сразу: их
объединение $F[(F^\times)^{1/n}]$ имеет проконечную группу Галуа, двойственную дискретной группе $F^\times/F^{\times n}$;
это изложено в \rf{7.37}.

\paragraph{Теория Куммера и арифметика.} Все доказательства этого раздела~--- чистая теория Галуа: основная теорема
(\rf{3.17}), теорема Дедекинда (\rf{5.14}), теорема Гильберта~90 (\rf{5.23}) и двойственность конечных абелевых групп. Арифметика
конкретного поля входит только через строение группы $F^\times/F^{\times n}$: для $\Q$~--- это однозначность разложения
на простые (пример \rf{5.31}(b)), для $\Q[\omega]$~--- норма (\rf{5.31c}), для $\F_q$~--- цикличность $\F_q^\times$ (\rf{5.31f}), для
$\Qp$~--- лемма Гензеля (\rf{5.31h}). Для числовых и локальных полей все абелевы расширения (и без условия $\mu_n\subset F$)
описывает теория полей классов; спаривание \rf{5.32} вместе с её отображением взаимности даёт символ Гильберта (о теории
полей классов см.\ \rf{7.39}).
